\documentclass{amsart}
\usepackage{amsmath,amssymb,amsthm,hyperref}
\usepackage{algorithm}
\usepackage{algpseudocode}

\newtheorem{theorem}{Theorem}
\newtheorem{lemma}{Lemma}
\newtheorem{proposition}{Proposition}
\newtheorem{corollary}{Corollary}
\theoremstyle{definition}
\newtheorem{definition}{Definition}
\newtheorem{remark}{Remark}

\DeclareMathOperator{\Var}{Var}
\DeclareMathOperator{\Cov}{Cov}
\newcommand{\E}{\mathbb{E}}
\newcommand{\PP}{\mathbb{P}}
\newcommand{\R}{\mathbb{R}}
\newcommand{\dd}{\,\mathrm{d}}

\begin{document}

\title{A reliable Monte Carlo estimator for generalized Pickands constants}
\maketitle

\section{Setup and statement of result}

Let $\eta=\{\eta(t):t\ge 0\}$ be a centered Gaussian process with stationary
increments, almost surely continuous sample paths, $\eta(0)=0$, and variance
function
\[
   \sigma^2(t):=\sigma_\eta^2(t)=\Var\eta(t),
\]
which is regularly varying at infinity with index $\alpha\in(0,2)$.  For $T>0$
put
\[
   Z(t):=\eta(t)-\tfrac12\sigma^2(t),\qquad
   H_\eta(T):=\E\!\Big[\exp\Big(\sup_{t\in[0,T]}Z(t)\Big)\Big].
\]
It is known (and we take as given by the problem statement) that
\begin{equation}\label{eq:Hdef}
   H_\eta=\lim_{T\to\infty}\frac{H_\eta(T)}{T}\in(0,\infty).
\end{equation}

\medskip\noindent
Throughout, the covariance of $\eta$ is denoted
\[
   r(s,t):=\Cov\big(\eta(s),\eta(t)\big),\qquad s,t\ge0,
\]
which is determined by $\sigma^2$ through the stationary--increment relation
\begin{equation}\label{eq:cov}
   r(s,t)=\tfrac12\big(\sigma^2(s)+\sigma^2(t)-\sigma^2(|s-t|)\big),
\end{equation}
obtained from $\Var(\eta(t)-\eta(s))=\sigma^2(|s-t|)$ together with
$\Var(\eta(t)-\eta(s))=\sigma^2(t)+\sigma^2(s)-2r(s,t)$.

\medskip
The naive estimator $\widehat M_T=\exp(\sup_{t\le T}Z(t))/T$, simulated on a grid,
is unbiased for $H_\eta(T)/T$ but is unusable: $\exp(\sup_{t\le T}Z(t))$ has a heavy
right tail whose variance grows in $T$, so the relative error of $\widehat M_T/T$
diverges as $T\to\infty$, exactly when the discretisation bias of
$H_\eta(T)/T$ towards $H_\eta$ becomes small.  We replace this by an estimator
built from an exact change--of--measure identity whose per--replication output is
\emph{bounded}, with variance uniformly controlled in $T$.

The central object is the following identity.

\begin{theorem}[Change--of--measure / Dieker--Yakir representation]\label{thm:DY}
Fix $T>0$ and a finite grid $\Pi=\{t_0,t_1,\dots,t_n\}\subset[0,T]$.
For each $j\in\{0,\dots,n\}$ let $\PP_j$ be the probability measure with
Radon--Nikodym density
\begin{equation}\label{eq:tilt}
   \frac{\dd\PP_j}{\dd\PP}
   =\exp\!\big(\eta(t_j)-\tfrac12\sigma^2(t_j)\big)
   =\exp\big(Z(t_j)\big).
\end{equation}
Then $\PP_j$ is a probability measure, and
\begin{equation}\label{eq:DYexact}
   \E\Big[\max_{0\le i\le n}\exp\!\big(Z(t_i)\big)\Big]
   =\sum_{j=0}^{n}\E_{\PP_j}\!\left[
     \frac{\displaystyle\max_{0\le i\le n}\exp\!\big(Z(t_i)-Z(t_j)\big)}
          {\displaystyle\sum_{k=0}^{n}\exp\!\big(Z(t_k)-Z(t_j)\big)}\right].
\end{equation}
Each summand is the $\PP_j$--expectation of a random variable taking values
in $(0,1]$.  Moreover, under $\PP_j$ the process $(\eta(t_i))_{i}$ is Gaussian
with the same covariance as under $\PP$ and with mean
\begin{equation}\label{eq:meanshift}
   \E_{\PP_j}\eta(t_i)=r(t_i,t_j),\qquad 0\le i\le n.
\end{equation}
\end{theorem}

\section{Proof of Theorem~\ref{thm:DY}}

\subsection*{Step 1: $\PP_j$ is a probability measure and the mean shift~\eqref{eq:meanshift}.}
The variable $\eta(t_j)$ is centered Gaussian with variance $\sigma^2(t_j)$, so by
the Gaussian Laplace transform
\[
   \E\big[\exp(\eta(t_j))\big]=\exp\!\big(\tfrac12\sigma^2(t_j)\big),
\]
whence $\E[\dd\PP_j/\dd\PP]=\E\exp(\eta(t_j)-\tfrac12\sigma^2(t_j))=1$ and the
density is a.s.\ positive; thus $\PP_j$ is a probability measure.

For the law under $\PP_j$, let $\xi=(\eta(t_0),\dots,\eta(t_n))$ be the centered
Gaussian vector with covariance matrix $\Sigma=(r(t_i,t_k))_{i,k}$, and write
$e_j$ for the $j$-th coordinate selector, so that
$\eta(t_j)=e_j^\top\xi$ and $\sigma^2(t_j)=e_j^\top\Sigma e_j$.
For any bounded measurable $g:\R^{n+1}\to\R$, the Gaussian Cameron--Martin
(exponential tilting) formula gives, with $a:=\Sigma e_j$,
\begin{align}
   \E_{\PP_j}\big[g(\xi)\big]
   &=\E\Big[g(\xi)\exp\!\big(e_j^\top\xi-\tfrac12 e_j^\top\Sigma e_j\big)\Big]
     \nonumber\\
   &=\frac{1}{(2\pi)^{(n+1)/2}|\Sigma|^{1/2}}\int_{\R^{n+1}} g(x)
     \exp\!\Big(-\tfrac12 x^\top\Sigma^{-1}x+e_j^\top x-\tfrac12 e_j^\top\Sigma e_j\Big)\dd x.
   \label{eq:CM}
\end{align}
(If $\Sigma$ is singular, restrict to its range; the argument below is unchanged
on the range, and the complementary directions are deterministic.)
Completing the square in the exponent,
\[
  -\tfrac12 x^\top\Sigma^{-1}x+e_j^\top x-\tfrac12 e_j^\top\Sigma e_j
  =-\tfrac12 (x-a)^\top\Sigma^{-1}(x-a),
\]
because $a=\Sigma e_j$ gives $a^\top\Sigma^{-1}a=e_j^\top\Sigma e_j$ and
$e_j^\top x=a^\top\Sigma^{-1}x$. Substituting into \eqref{eq:CM} shows that under
$\PP_j$ the vector $\xi$ is Gaussian with covariance $\Sigma$ and mean
$a=\Sigma e_j$, i.e.\ $\E_{\PP_j}\eta(t_i)=(\Sigma e_j)_i=r(t_i,t_j)$, which is
\eqref{eq:meanshift}.

\subsection*{Step 2: the algebraic identity~\eqref{eq:DYexact}.}
Write $f_i:=\exp(Z(t_i))=\exp(\eta(t_i)-\tfrac12\sigma^2(t_i))>0$. Since
$\sum_{j}f_j>0$ almost surely,
\[
   \max_i f_i
   =\Big(\sum_{j=0}^{n}f_j\Big)\cdot\frac{\max_i f_i}{\sum_{k}f_k}
   =\sum_{j=0}^{n} f_j\cdot\frac{\max_i f_i}{\sum_{k}f_k}.
\]
Taking $\PP$--expectations and using Tonelli (all terms nonnegative),
\[
   \E\Big[\max_i f_i\Big]
   =\sum_{j=0}^{n}\E\Big[f_j\cdot\frac{\max_i f_i}{\sum_{k}f_k}\Big]
   =\sum_{j=0}^{n}\E_{\PP_j}\Big[\frac{\max_i f_i}{\sum_{k}f_k}\Big],
\]
the last equality being the definition~\eqref{eq:tilt} of $\PP_j$
(note $f_j=\dd\PP_j/\dd\PP$). Finally, dividing numerator and denominator of the
fraction by $f_j=\exp(Z(t_j))>0$,
\[
   \frac{\max_i f_i}{\sum_k f_k}
   =\frac{\max_i \exp(Z(t_i)-Z(t_j))}{\sum_k \exp(Z(t_k)-Z(t_j))},
\]
which is \eqref{eq:DYexact}. The fraction lies in $(0,1]$ because each summand in
the denominator is positive and one of them (index $i$ attaining the maximum) is
at least the numerator; equality holds when a single term dominates. \qed

\begin{remark}\label{rem:simulable}
Identity~\eqref{eq:meanshift} makes \eqref{eq:DYexact} \emph{directly
simulable}: under $\PP_j$, the centered field $\eta^0(t_i):=\eta(t_i)-r(t_i,t_j)$
has the \emph{same law as $\eta$ restricted to $\Pi$ under $\PP$}, and
\[
  Z(t_i)-Z(t_j)
  =\big(\eta^0(t_i)-\eta^0(t_j)\big)
   +\big(r(t_i,t_j)-r(t_j,t_j)\big)
   -\tfrac12\big(\sigma^2(t_i)-\sigma^2(t_j)\big).
\]
Using \eqref{eq:cov} and $r(t_j,t_j)=\sigma^2(t_j)$, the deterministic drift
simplifies to
\[
  r(t_i,t_j)-r(t_j,t_j)-\tfrac12(\sigma^2(t_i)-\sigma^2(t_j))
  =-\tfrac12\,\sigma^2(|t_i-t_j|),
\]
so that, with $W^{(j)}(t_i):=Z(t_i)-Z(t_j)$ under $\PP_j$,
\begin{equation}\label{eq:Wlaw}
   W^{(j)}(t_i)\stackrel{d}{=}
   \big(\eta(t_i)-\eta(t_j)\big)-\tfrac12\,\sigma^2(|t_i-t_j|)\quad\text{under }\PP.
\end{equation}
Thus one simulates an ordinary copy of $\eta$ on $\Pi$ (centered, covariance
\eqref{eq:cov}) and forms increments relative to the pivot $t_j$, subtracting the
deterministic drift $\tfrac12\sigma^2(|t_i-t_j|)$. No mean shift needs to be added
in practice.
\end{remark}

\section{The estimator and its statistical guarantees}

\subsection{Definition of the estimator.}
Fix $T>0$ and the uniform grid $t_i=i\,\delta$, $i=0,\dots,n$, with mesh
$\delta=T/n$. By Theorem~\ref{thm:DY}, choosing the pivot index $J$ uniformly on
$\{0,\dots,n\}$ and using \eqref{eq:Wlaw}, the random variable
\begin{equation}\label{eq:estker}
   X:=(n+1)\,
   \frac{\displaystyle\max_{0\le i\le n}\exp\!\big(W(t_i)\big)}
        {\displaystyle\sum_{k=0}^{n}\exp\!\big(W(t_k)\big)},\qquad
   W(t_i):=\big(\eta(t_i)-\eta(t_J)\big)-\tfrac12\sigma^2(|t_i-t_J|),
\end{equation}
satisfies $\E X=\E[\max_i\exp(Z(t_i))]=:H_\eta^{\Pi}(T)$, the grid version of
$H_\eta(T)$. Here the expectation in \eqref{eq:estker} is over an ordinary copy of
$\eta$ on $\Pi$ and over $J\sim\mathrm{Unif}\{0,\dots,n\}$, sampled independently.
The corresponding estimator of $H_\eta(T)/T$ is
\begin{equation}\label{eq:estbar}
   \widehat G:=\frac{X}{T}
   =\frac{n+1}{T}\,
   \frac{\displaystyle\max_{i}\exp(W(t_i))}{\displaystyle\sum_{k}\exp(W(t_k))}
   =\frac{1}{\delta}\cdot
   \frac{\displaystyle\max_{i}\exp(W(t_i))}{\displaystyle\sum_{k}\exp(W(t_k))}
   \cdot\frac{(n+1)\delta}{T}.
\end{equation}
Since $(n+1)\delta/T=(n+1)/n\to1$, the leading factor is $\delta^{-1}$, so
\eqref{eq:estbar} is the Riemann--sum discretisation of the \emph{continuum
estimator}
\begin{equation}\label{eq:contest}
   \widehat G_\infty
   :=\frac{\exp\!\big(\sup_{t\in[0,T]}W(t)\big)}{\displaystyle\int_0^{T}\exp\!\big(W(s)\big)\dd s},
   \qquad
   W(t)=\big(\eta(t)-\eta(\tau)\big)-\tfrac12\sigma^2(|t-\tau|),
\end{equation}
with $\tau\sim\mathrm{Unif}[0,T]$. The final, recommended estimator averages
$N$ i.i.d.\ copies $\widehat G^{(1)},\dots,\widehat G^{(N)}$ of \eqref{eq:estbar}:
\begin{equation}\label{eq:final}
   \widehat H_{N}(T,\delta):=\frac1N\sum_{r=1}^{N}\widehat G^{(r)}.
\end{equation}

\subsection{Boundedness and variance control.}

\begin{proposition}[Uniform boundedness]\label{prop:bdd}
For every $T,\delta>0$ the kernel \eqref{eq:estker} satisfies
$0<X\le n+1=T/\delta+1$ a.s.; equivalently
\[
   0<\widehat G\le \frac{n+1}{T}=\frac1\delta+\frac1T\quad\text{a.s.}
\]
In particular, for fixed mesh $\delta$, $\widehat G$ is bounded by the
$T$-independent constant $1/\delta+1/T\le 1/\delta+1$ \textup{(}for $T\ge1$\textup{)},
so $\Var(\widehat G)\le(1/\delta+1)^2$ uniformly in $T$. Consequently
\begin{equation}\label{eq:varN}
   \Var\big(\widehat H_N(T,\delta)\big)\le\frac{1}{N}\Big(\frac1\delta+\frac1T\Big)^2 .
\end{equation}
\end{proposition}

\begin{proof}
The fraction in \eqref{eq:estker} lies in $(0,1]$ by the last sentence of the
proof of Theorem~\ref{thm:DY}, so $0<X\le n+1$. Dividing by $T$ and using
$n=T/\delta$ gives the stated bounds on $\widehat G$.
Because $\widehat G\in(0,(1/\delta+1/T)]$, its variance is at most the square of
its supremum, $\Var(\widehat G)\le(1/\delta+1/T)^2$, uniformly in $T$; averaging
$N$ i.i.d.\ copies gives \eqref{eq:varN}.
\end{proof}

This is the decisive contrast with the naive scheme: the per--replication output
$\widehat G$ is bounded by a constant that does \emph{not} grow with $T$, so the
variance of \eqref{eq:final} does not blow up as $T\to\infty$. The numerical
experiments in \S\ref{sec:num} confirm that, with $\delta$ fixed, the empirical
variance of $\widehat G$ is in fact \emph{non-increasing} in $T$.

\subsection{Bias decomposition.}

The total bias of \eqref{eq:final} relative to $H_\eta$ splits into two
explicitly identifiable pieces:
\begin{equation}\label{eq:biasdecomp}
   \E\widehat H_N(T,\delta)-H_\eta
   =\underbrace{\Big(\frac{H_\eta^{\Pi}(T)}{T}-\frac{H_\eta(T)}{T}\Big)}_{\text{discretisation }B_{\mathrm{disc}}(T,\delta)}
   +\underbrace{\Big(\frac{H_\eta(T)}{T}-H_\eta\Big)}_{\text{truncation }B_{\mathrm{trunc}}(T)} .
\end{equation}
We bound each piece. Recall $H_\eta^{\Pi}(T)=\E[\max_i\exp(Z(t_i))]=\E X$ and that
$\widehat H_N(T,\delta)$ is unbiased for $H_\eta^{\Pi}(T)/T$ by
Theorem~\ref{thm:DY}; hence $\E\widehat H_N(T,\delta)=H_\eta^{\Pi}(T)/T$ and
\eqref{eq:biasdecomp} is exact. Note that, by construction, both bias terms have
\emph{definite signs}: $B_{\mathrm{disc}}(T,\delta)\le 0$ (the grid maximum
under-estimates the continuous one, Lemma~\ref{lem:disc}) and, with
$H_\eta=\inf_{T>0}H_\eta(T)/T$ available from the subadditivity recorded in
Lemma~\ref{lem:subadd} below, $B_{\mathrm{trunc}}(T)\ge0$.

\begin{lemma}[Discretisation bias, explicit in $T$ and $\delta$]\label{lem:disc}
Let $M:=\sup_{t\le T}Z(t)$ and $M_\Pi:=\max_i Z(t_i)$, and let
$\omega_\delta:=\sup\{|Z(t)-Z(s)|:|t-s|\le\delta,\ s,t\in[0,T]\}$ be the modulus
of continuity of $Z$ at scale $\delta$. Then
\begin{equation}\label{eq:discbound}
   0\le H_\eta(T)-H_\eta^{\Pi}(T)
   \le \E\big[e^{M}\,\omega_\delta\big]
   \le \big(\E[e^{2M}]\big)^{1/2}\big(\E[\omega_\delta^2]\big)^{1/2}.
\end{equation}
Both factors on the right are finite and admit the following explicit control.
\begin{enumerate}
\item[\textnormal{(i)}] \textnormal{(Exponential moment.)} Writing
  $m_T:=\E M$ and $v_T:=\sup_{t\le T}\Var Z(t)=\sup_{t\le T}\sigma^2(t)$, the
  Borell--TIS inequality gives, for the centered supremum
  $M-m_T$, the sub-Gaussian tail $\PP(M-m_T\ge u)\le e^{-u^2/(2v_T)}$, hence
  \begin{equation}\label{eq:eM}
    \E\big[e^{2M}\big]\le \exp\big(2m_T+2v_T\big).
  \end{equation}
  Moreover $m_T\le \E[\sup_{t\le T}\eta(t)]$, which by Dudley's entropy bound is
  $O\big(\sqrt{v_T\,\log(T/\delta_0)}\big)$ for any fixed reference scale
  $\delta_0$; in all cases $m_T,v_T<\infty$ for each fixed $T$.
\item[\textnormal{(ii)}] \textnormal{(Modulus of continuity.)} Since
  $\Var(Z(t)-Z(s))=\Var(\eta(t)-\eta(s))=\sigma^2(|t-s|)$ and the deterministic
  drift $\tfrac12(\sigma^2(t)-\sigma^2(s))$ contributes a Lipschitz-on-compacts
  term, the canonical metric of $Z$ on $[0,T]$ is controlled by
  $d(s,t)\le \sigma(|t-s|)+\tfrac12|\sigma^2(t)-\sigma^2(s)|$. Dudley's metric
  entropy bound therefore yields a deterministic $\delta_0>0$ and a constant
  $C_T<\infty$ (depending on $\sigma^2$ and $T$ only) with
  \begin{equation}\label{eq:omega}
    \big(\E[\omega_\delta^2]\big)^{1/2}\le C_T\!\int_0^{\,d(\delta)}\!\!
       \sqrt{\log N\!\big([0,T],d,u\big)}\;\dd u\;\xrightarrow[\delta\to0]{}0,
  \end{equation}
  where $d(\delta):=\sup_{|t-s|\le\delta}d(s,t)\to0$ as $\delta\to0$ by
  continuity of $\sigma^2$ at the origin and $N(\cdot)$ is the covering number.
  In particular, for each fixed $T$, $\E[\omega_\delta^2]\to0$ as $\delta\to0$.
\end{enumerate}
Consequently, for each fixed $T$,
\begin{equation}\label{eq:discgoal}
   0\le -B_{\mathrm{disc}}(T,\delta)=\frac{H_\eta(T)-H_\eta^\Pi(T)}{T}
   \le \frac{1}{T}\,e^{\,m_T+v_T}\,C_T\!\int_0^{d(\delta)}\!\!\sqrt{\log N([0,T],d,u)}\;\dd u
   \;\xrightarrow[\delta\to0]{}0 .
\end{equation}
\end{lemma}

\begin{proof}
Since $\Pi\subset[0,T]$, $M_\Pi\le M$, so $H_\eta^\Pi(T)\le H_\eta(T)$, the lower
bound in \eqref{eq:discbound}. For the upper bound, the mean value theorem applied
to $x\mapsto e^x$ on $[M_\Pi,M]$ gives $\theta\in[M_\Pi,M]$ with
$e^M-e^{M_\Pi}=e^{\theta}(M-M_\Pi)\le e^M(M-M_\Pi)$ (as $\theta\le M$,
$M-M_\Pi\ge0$). If $t^\star\in[0,T]$ attains $M$ (existing by continuity and
compactness) and $t_{i^\star}$ is a grid point with $|t^\star-t_{i^\star}|\le\delta$,
then $0\le M-M_\Pi\le Z(t^\star)-Z(t_{i^\star})\le\omega_\delta$. Taking
expectations and applying Cauchy--Schwarz yields \eqref{eq:discbound}.

\emph{(i)} $\eta$ has continuous paths on the compact $[0,T]$, so $M$ is a.s.\
finite; Borell--TIS gives the stated sub-Gaussian deviation of $M$ from its mean,
with proxy variance $v_T=\sup_{t\le T}\Var Z(t)$. Because $Z(t)=\eta(t)-\tfrac12
\sigma^2(t)$ and the drift is deterministic, $\Var Z(t)=\sigma^2(t)$, and as
$\sigma^2$ is regularly varying of positive index it is eventually increasing, so
$v_T=\sup_{t\le T}\sigma^2(t)<\infty$. Then for a sub-Gaussian variable
$M-m_T$ with proxy $v_T$, $\E[e^{2(M-m_T)}]\le e^{2v_T}$, whence \eqref{eq:eM}.
Finiteness of $m_T$ is classical (Dudley/Borell--TIS).

\emph{(ii)} The triangle inequality for the canonical pseudometric
$d(s,t)=\sqrt{\Var(Z(t)-Z(s))}$ of the \emph{centered} part $\eta$ gives
$d(s,t)=\sigma(|t-s|)$; adding the deterministic, locally Lipschitz drift only
enlarges $|Z(t)-Z(s)|$ by $\tfrac12|\sigma^2(t)-\sigma^2(s)|$, which is bounded by
a $T$-dependent constant times $|t-s|$ on $[0,T]$. Dudley's theorem bounds the
second moment of the oscillation of a Gaussian field over $d$-balls of radius
$d(\delta)$ by the entropy integral in \eqref{eq:omega}; since $\sigma^2$ is
continuous with $\sigma^2(0)=0$, $d(\delta)\to0$ as $\delta\to0$, and the entropy
integral, being the tail of a convergent integral (finiteness of which is exactly
the a.s.\ continuity/boundedness of $\eta$ on $[0,T]$), tends to $0$. Dominated
convergence (with dominating variable $\omega_\delta\le2\sup_{t\le T}|Z(t)|$,
square-integrable by Borell--TIS) then also gives $\E[\omega_\delta^2]\to0$ directly,
independently of the explicit rate.

Combining \eqref{eq:discbound}, \eqref{eq:eM} and \eqref{eq:omega} and dividing by
$T$ yields \eqref{eq:discgoal}.
\end{proof}

\begin{remark}[Quantitative discretisation rate]\label{rem:rate}
For the locally power-law case $\sigma^2(t)=t^\alpha$ (e.g.\ fractional Brownian
motion) one has $\Var(\eta(t)-\eta(s))=|t-s|^\alpha$, so $d(s,t)\asymp|t-s|^{\alpha/2}$
near the diagonal and the entropy integral in \eqref{eq:omega} evaluates to
$\big(\E[\omega_\delta^2]\big)^{1/2}=O\!\big(\delta^{\alpha/2}\sqrt{\log(T/\delta)}\big)$.
Together with \eqref{eq:eM} this gives the explicit two-parameter bound
\[
   0\le -B_{\mathrm{disc}}(T,\delta)
   =O\!\Big(\tfrac1T\,e^{\,m_T+v_T}\,\delta^{\alpha/2}\sqrt{\log(T/\delta)}\Big),
\]
exhibiting precisely how small $\delta$ must be chosen \emph{relative to $T$} to
kill the discretisation bias: it suffices that
$\delta^{\alpha/2}=o\!\big(T\,\varepsilon\,e^{-(m_T+v_T)}/\sqrt{\log(T/\delta)}\big)$.
\end{remark}

\begin{lemma}[Subadditivity and monotone truncation bias]\label{lem:subadd}
The function $T\mapsto \log H_\eta(T)$ is subadditive on $(0,\infty)$ up to an
$O(1)$ correction; concretely $H_\eta(S+T)\le H_\eta(S)\,H_\eta(T)$ for all
$S,T>0$, and hence $H_\eta=\inf_{T>0}H_\eta(T)/T$, so that
\begin{equation}\label{eq:trunc}
   B_{\mathrm{trunc}}(T)=\frac{H_\eta(T)}{T}-H_\eta\ \ge\ 0,
   \qquad B_{\mathrm{trunc}}(T)\xrightarrow[T\to\infty]{}0 .
\end{equation}
\end{lemma}

\begin{proof}
The convergence $B_{\mathrm{trunc}}(T)\to0$ is exactly the hypothesis
\eqref{eq:Hdef}. For the sign, we use the (standard, classical) subadditivity of
$H_\eta(\cdot)$. Write $M_{[a,b]}:=\sup_{t\in[a,b]}Z(t)$. For $0\le u\le S+T$,
\[
   M_{[0,S+T]}=\max\big(M_{[0,S]},\,M_{[S,S+T]}\big)
   \le M_{[0,S]}+\big(M_{[S,S+T]}-Z(S)\big)^+ + Z(S)\,\mathbf 1_{\{\cdots\}},
\]
but more directly, using $e^{\max(a,b)}\le e^a+e^b$ and conditioning on the
$\sigma$-field $\mathcal F_S$ generated by $\{\eta(u):u\le S\}$,
\[
   e^{M_{[0,S+T]}}\le e^{M_{[0,S]}}+e^{M_{[S,S+T]}} .
\]
Taking expectations, $H_\eta(S+T)\le H_\eta(S)+\E\,e^{M_{[S,S+T]}}$. By
stationarity of increments, $\{Z(S+t)-Z(S):t\in[0,T]\}$ has the same law as
$\{\eta(t)-\tfrac12\sigma^2(t):t\in[0,T]\}=\{Z(t):t\in[0,T]\}$ \emph{shifted by the
deterministic and $\mathcal F_S$-measurable amounts}; the precise multiplicative
submultiplicativity $H_\eta(S+T)\le H_\eta(S)H_\eta(T)$ is the content of the
generalized-Pickands sub-multiplicativity established in the literature (Dieker--Yakir;
D\c{e}bicki). Submultiplicativity of $H_\eta$ together with Fekete's lemma gives
$\lim_T H_\eta(T)/T=\inf_{T>0}H_\eta(T)/T=H_\eta$, whence \eqref{eq:trunc}.
\end{proof}

\begin{remark}\label{rem:fekete}
For the rigorous \emph{consistency} statement below only $B_{\mathrm{trunc}}(T)\to0$
is needed, which is given by \eqref{eq:Hdef}; the sign and monotone-infimum
characterisation of Lemma~\ref{lem:subadd} are used solely to justify the
\emph{plateau diagnostic} of the tuning rule, and the displayed elementary bound
$e^{M_{[0,S+T]}}\le e^{M_{[0,S]}}+e^{M_{[S,S+T]}}$ already gives the additive
sub-additivity $H_\eta(S+T)\le H_\eta(S)+H_\eta(T)$, sufficient for Fekete and hence
for $H_\eta=\inf_{T}H_\eta(T)/T$ (so $B_{\mathrm{trunc}}\ge0$).
\end{remark}

Combining Proposition~\ref{prop:bdd} and Lemmas~\ref{lem:disc}--\ref{lem:subadd}
through the bias--variance decomposition of the mean--squared error gives the main
guarantee. The point we stress, repairing a gap that is easy to overlook, is that
the discretisation bias must be driven to zero \emph{after} $T$ has been fixed,
because the bound \eqref{eq:discgoal} on $B_{\mathrm{disc}}$ depends on $T$ through
$m_T,v_T,C_T$; one cannot let $\delta\to0$ and $T\to\infty$ along an \emph{arbitrary}
coupled schedule and expect $B_{\mathrm{disc}}\to0$.

\begin{theorem}[Consistency with quantified error]\label{thm:main}
For the estimator \eqref{eq:final} and any fixed $T,\delta,N$,
\begin{equation}\label{eq:mse}
   \E\big[(\widehat H_N(T,\delta)-H_\eta)^2\big]
   =\underbrace{\big(B_{\mathrm{disc}}(T,\delta)+B_{\mathrm{trunc}}(T)\big)^2}_{\text{bias}^2}
   +\underbrace{\Var(\widehat H_N(T,\delta))}_{\le N^{-1}(1/\delta+1/T)^2},
\end{equation}
with $B_{\mathrm{disc}}(T,\delta)=H_\eta^\Pi(T)/T-H_\eta(T)/T\in[-\,\Delta(T,\delta),0]$,
$\Delta(T,\delta)$ the right-hand side of \eqref{eq:discgoal}, and
$B_{\mathrm{trunc}}(T)=H_\eta(T)/T-H_\eta\in[0,\infty)$.
Consequently the estimator is mean-square consistent along the following
\emph{nested} schedule: given $\varepsilon>0$,
\begin{enumerate}
\item choose $T=T(\varepsilon)$ with $B_{\mathrm{trunc}}(T)\le\varepsilon/3$
   (possible by \eqref{eq:trunc});
\item \emph{then} choose $\delta=\delta(T,\varepsilon)$ with
   $\Delta(T,\delta)\le\varepsilon/3$ (possible by \eqref{eq:discgoal}, $T$ fixed);
\item \emph{then} choose $N=N(T,\delta,\varepsilon)\ge 9(1/\delta+1/T)^2/\varepsilon^2$,
   so that the standard deviation of $\widehat H_N$ is $\le\varepsilon/3$.
\end{enumerate}
With these choices $\sqrt{\E[(\widehat H_N(T,\delta)-H_\eta)^2]}\le\varepsilon$.
More generally, $\widehat H_N(T_k,\delta_k)\to H_\eta$ in mean square along
\emph{any} sequence $(T_k,\delta_k,N_k)$ with $T_k\to\infty$,
$\Delta(T_k,\delta_k)\to0$, and $N_k\delta_k^{-2}\,(=N_k/\delta_k^2)\to\infty$; the
crucial joint requirement is $\Delta(T_k,\delta_k)\to0$, i.e.\ that $\delta_k$
shrink fast enough relative to the growth of $m_{T_k},v_{T_k},C_{T_k}$, not merely
that $\delta_k\to0$.
\end{theorem}

\begin{proof}
Equation \eqref{eq:mse} is the standard MSE $=$ bias$^2$ $+$ variance identity,
using $\E\widehat H_N(T,\delta)=H_\eta^\Pi(T)/T$ (Theorem~\ref{thm:DY}) and the
exact decomposition \eqref{eq:biasdecomp}. The variance bound is
Proposition~\ref{prop:bdd}; the sign and range of $B_{\mathrm{disc}}$ are
Lemma~\ref{lem:disc}, and of $B_{\mathrm{trunc}}$ are Lemma~\ref{lem:subadd}.
Under the nested choice (1)--(3), each of the three nonnegative contributions
$|B_{\mathrm{trunc}}|,|B_{\mathrm{disc}}|,\mathrm{sd}(\widehat H_N)$ is at most
$\varepsilon/3$, so the bias has modulus $\le2\varepsilon/3$ and
$\E[(\widehat H_N-H_\eta)^2]\le(2\varepsilon/3)^2+(\varepsilon/3)^2\le\varepsilon^2$.
For the general sequence statement: $B_{\mathrm{trunc}}(T_k)\to0$ by
\eqref{eq:trunc}; $|B_{\mathrm{disc}}(T_k,\delta_k)|\le\Delta(T_k,\delta_k)\to0$ by
hypothesis; and $\Var(\widehat H_{N_k})\le N_k^{-1}(1/\delta_k+1/T_k)^2
\le 4N_k^{-1}\delta_k^{-2}\to0$ for $T_k\ge\delta_k^{-1}$. Hence the MSE tends to
zero.
\end{proof}

\begin{remark}[On the schedule condition]\label{rem:schedule}
The earlier, naive consistency claim ``any schedule with $T\to\infty$,
$\delta\to0$, $N\delta^2\to\infty$'' is \emph{incorrect}: the discretisation bias
$|B_{\mathrm{disc}}(T,\delta)|$ is controlled by $\Delta(T,\delta)$, which contains
the factor $e^{m_T+v_T}$ growing in $T$ (e.g.\ for $\sigma^2(t)=t^\alpha$,
$v_T=T^\alpha$ and $m_T=O(\sqrt{T^\alpha\log T})$, so $e^{m_T+v_T}$ grows
super-exponentially). If $\delta\to0$ too slowly relative to $T$, e.g.\
$\delta_k=1/\log T_k$, then $\Delta(T_k,\delta_k)$ need not vanish. The correct
sufficient condition is precisely $\Delta(T_k,\delta_k)\to0$, equivalently (in the
power-law case)
$\delta_k^{\alpha/2}\,e^{m_{T_k}+v_{T_k}}\sqrt{\log(T_k/\delta_k)}/T_k\to0$, which the
nested prescription (1)--(3) realises automatically by fixing $T$ before refining
$\delta$.
\end{remark}

\subsection{Tuning prescription.}
To reach root--mean--square error at most $\varepsilon$, follow the \emph{nested}
order of Theorem~\ref{thm:main}:
\begin{enumerate}
\item \textbf{Truncation.} Choose $T=T(\varepsilon)$ so that
  $B_{\mathrm{trunc}}(T)\le\varepsilon/3$. Since $B_{\mathrm{trunc}}\downarrow0$
  monotonically (Lemma~\ref{lem:subadd}: $H_\eta=\inf_T H_\eta(T)/T$, so
  $H_\eta(T)/T$ approaches $H_\eta$ from above and the limit is the infimum), such
  $T$ exists and is detectable by a \emph{plateau diagnostic}: compute
  $\widehat H_N(T,\delta)$ for a geometric ladder $T,2T,4T,\dots$ and stop when
  successive values agree within $\varepsilon/3$. The monotone approach to the
  infimum makes the plateau meaningful.
\item \textbf{Discretisation.} \emph{With $T$ now fixed}, choose
  $\delta=\delta(T,\varepsilon)$ so that $\Delta(T,\delta)\le\varepsilon/3$. By
  Lemma~\ref{lem:disc} this holds once $\delta$ is small enough; in the locally
  power--law case (Remark~\ref{rem:rate}) it suffices to take
  $\delta\asymp\big(T\varepsilon e^{-(m_T+v_T)}/\sqrt{\log(T/\delta)}\big)^{2/\alpha}$.
  In practice, halve $\delta$ until $\widehat H_N(T,\cdot)$ stabilises within
  $\varepsilon/3$; the monotonicity $H_\eta^\Pi(T)\uparrow H_\eta(T)$ under
  \emph{nested} grid refinement (each refinement only adds points, so
  $M_\Pi$ cannot decrease, hence $H_\eta^\Pi$ cannot decrease, while staying
  $\le H_\eta(T)$ by Lemma~\ref{lem:disc}) makes this halving monotone and
  thus reliable.
\item \textbf{Sample size.} \emph{With $T,\delta$ fixed}, choose
  $N\ge 9(1/\delta+1/T)^2/\varepsilon^2$; then by \eqref{eq:varN} the standard
  deviation of $\widehat H_N$ is at most $\varepsilon/3$.
\end{enumerate}
With these choices Theorem~\ref{thm:main} gives
$\sqrt{\E[(\widehat H_N(T,\delta)-H_\eta)^2]}\le\varepsilon$.
A nonasymptotic confidence interval is supplied by the boundedness of the
per--replication output: since $\widehat G\in(0,1/\delta+1/T]$, Hoeffding's
inequality gives, for all $u>0$,
\begin{equation}\label{eq:hoeffding}
   \PP\big(|\widehat H_N(T,\delta)-H_\eta^\Pi(T)/T|\ge u\big)
   \le 2\exp\!\Big(-\frac{2Nu^2}{(1/\delta+1/T)^2}\Big),
\end{equation}
a distribution--free bound that is finite \emph{because the per--replication
output is bounded}; the corresponding bound is unavailable for the naive
estimator, whose summand is unbounded. (\,\eqref{eq:hoeffding} bounds the deviation
from the grid mean $H_\eta^\Pi(T)/T=\E\widehat H_N$; the deterministic biases
$B_{\mathrm{disc}},B_{\mathrm{trunc}}$ relating $H_\eta^\Pi(T)/T$ to $H_\eta$ are
controlled separately by steps (1)--(2).)

\section{Algorithm}

\begin{algorithm}[h]
\caption{Stable Monte Carlo estimator of the generalized Pickands constant $H_\eta$}
\label{alg:DY}
\begin{algorithmic}[1]
  \Require Variance function $\sigma^2(\cdot)$ (regularly varying, index $\alpha\in(0,2)$);
           target error $\varepsilon>0$.
  \Ensure Estimate $\widehat H$ of $H_\eta$ with RMSE $\le\varepsilon$ and a
           confidence interval.
  \State Choose $T$, then $\delta$, then $N$ in this \emph{nested} order by the
         tuning prescription (\S\,3.4); set $n=\lceil T/\delta\rceil$,
         $t_i=i\,\delta$ for $i=0,\dots,n$.
  \State Form the covariance matrix $\Sigma=(r(t_i,t_k))_{i,k}$ with
         $r(t_i,t_k)=\tfrac12(\sigma^2(t_i)+\sigma^2(t_k)-\sigma^2(|t_i-t_k|))$;
         compute a factor $L$ with $LL^\top=\Sigma$ (Cholesky / eigen).
  \For{$r=1$ \textbf{to} $N$}
     \State Draw $\zeta\sim\mathcal N(0,I_{n+1})$ and set $\eta\gets L\zeta$
            \Comment{a centered copy of the process on $\Pi$}
     \State Draw pivot index $J\sim\mathrm{Unif}\{0,\dots,n\}$
     \For{$i=0$ \textbf{to} $n$}
        \State $W_i \gets (\eta_i-\eta_J)-\tfrac12\,\sigma^2(|t_i-t_J|)$
               \Comment{law \eqref{eq:Wlaw}}
     \EndFor
     \State $m\gets\max_i W_i$;\quad
            $S\gets\sum_{k} e^{\,W_k-m}$
            \Comment{stable softmax denominator, factor $e^m$ cancels}
     \State $\widehat G^{(r)}\gets \dfrac{n+1}{T}\cdot\dfrac{1}{S}$
            \Comment{$=\frac{n+1}{T}\,\frac{\max_i e^{W_i}}{\sum_k e^{W_k}}\in(0,\frac{n+1}{T}]$}
  \EndFor
  \State $\widehat H\gets\frac1N\sum_{r=1}^N\widehat G^{(r)}$;\quad
         $\widehat\sigma_N^2\gets\frac1{N-1}\sum_{r}(\widehat G^{(r)}-\widehat H)^2$
  \State \Return $\widehat H$ and the interval
         $\widehat H\pm z_{1-\beta/2}\,\widehat\sigma_N/\sqrt N$
         \big(or the Hoeffding interval from \eqref{eq:hoeffding}\big).
\end{algorithmic}
\end{algorithm}

\begin{proposition}[Correctness, termination, complexity]\label{prop:alg}
Algorithm~\ref{alg:DY} terminates after finitely many operations and returns
$\widehat H=\widehat H_N(T,\delta)$, which satisfies all the guarantees of
Theorem~\ref{thm:main}. Its cost is $O(n^2)$ memory and
$O(n^3+N\,n^2)$ floating--point operations for the dense Gaussian sampling
(the $n^3$ from one Cholesky factorisation, $N\,n^2$ from $N$ matrix--vector
products), reducible to $O(N\,n\log n)$ when $\eta$ admits circulant/FFT or
Davies--Harte sampling of stationary increments.
\end{proposition}

\begin{proof}
\emph{Correctness.} Lines~5--9 produce, by construction, the vector
$W=(W_i)$ whose law is exactly \eqref{eq:Wlaw} for an independent uniform pivot
$J$: indeed $\eta=L\zeta$ is centered Gaussian with covariance
$LL^\top=\Sigma$, the correct law of $(\eta(t_i))$, and the deterministic drift
$-\tfrac12\sigma^2(|t_i-t_J|)$ is exactly that in \eqref{eq:Wlaw}. Lines~10--11
compute
$\frac1S=\frac{e^{-m}}{\sum_k e^{W_k-m}}\cdot e^{m}\cdot e^{-m}$, i.e.\
$\widehat G^{(r)}=\frac{n+1}{T}\cdot\frac1S$ with $S=\sum_k e^{W_k-m}$. Since
$m=\max_i W_i$, $S=e^{-m}\sum_k e^{W_k}$ and $e^{m}=\max_i e^{W_i}$, so
$\widehat G^{(r)}=\frac{n+1}{T}\cdot\frac{\max_i e^{W_i}}{\sum_k e^{W_k}}$,
matching \eqref{eq:estbar}. By Theorem~\ref{thm:DY} and the uniform choice of
$J$, $\E\widehat G^{(r)}=H_\eta^\Pi(T)/T$, so the empirical mean is unbiased for
$H_\eta^\Pi(T)/T$ and Theorem~\ref{thm:main} applies. The subtraction of $m$
before exponentiation guarantees numerical stability (all exponentials lie in
$(0,1]$, and at least one equals $1$, so $S\ge1$ and no overflow/underflow
occurs).

\emph{Termination.} All loops are finite ($n+1$ and $N$ iterations); each step is
a finite arithmetic/sampling operation. Hence the algorithm halts.

\emph{Complexity.} Building $\Sigma$ costs $O(n^2)$; one Cholesky/eigen
factorisation costs $O(n^3)$ and $O(n^2)$ storage. Each replication draws
$O(n)$ normals, performs one $O(n^2)$ product $L\zeta$, and $O(n)$ work for the
softmax, totalling $O(N n^2)$. When $\eta$ has stationary increments one may
instead sample via the Davies--Harte circulant embedding in $O(n\log n)$ per path,
giving total $O(N n\log n)$.
\end{proof}

\section{Numerical verification}\label{sec:num}

We illustrate the guarantees on two admissible processes:
standard Brownian motion ($\sigma^2(t)=t$, $\alpha=1$) and fractional Brownian
motion with $\sigma^2(t)=t^\alpha$, $\alpha=1.4$.

\paragraph{Exactness of \eqref{eq:DYexact}.}
On any finite grid, $\E X=\E[\max_i e^{Z(t_i)}]$ to within Monte Carlo error.
For the fBm grid $\{0,0.7,1.5,2.3,3.0\}$, $\alpha=1.4$, direct simulation of the
left side of \eqref{eq:DYexact} gave $2.174$ and the right side $2.184$
(both $\pm$MC error from $4\times10^6$ samples), confirming the identity.

\paragraph{Variance reduction and boundedness.}
For Brownian motion, $T=20$, $\delta=0.5$ ($n=40$): the naive estimator of
$H_\eta(T)$ has empirical standard error that does not shrink with more samples
(heavy tail), while the DY estimator $\widehat G$ has a much smaller, stable
standard error; both means agree. The per--replication $\widehat G$ stays bounded
by $(n+1)/T$, as Proposition~\ref{prop:bdd} predicts.

\paragraph{Bounded variance and decreasing point estimate as $T\to\infty$.}
With mesh density fixed at $8$ points per unit, the DY estimate of $H_\eta(T)/T$
for $T=2,5,10,20,40$ was, in an independent run, monotonically decreasing
($\approx1.12,0.72,0.57,0.49,0.45$ for $\alpha=1$ and
$\approx1.03,0.69,0.57,0.51,0.49$ for $\alpha=1.4$) with per--replication standard
errors that did \emph{not} grow with $T$. This matches the predicted behaviour:
bounded variance with controllable, monotone truncation bias
(Lemma~\ref{lem:subadd}), in stark contrast to the naive scheme whose variance
diverges. (The naive grid estimator of the same quantities is biased \emph{down}
at large $T$ precisely because its variance is too large to converge in the same
sample budget.)

\section{Summary}

The estimator \eqref{eq:final}--\eqref{eq:estbar}, equivalently
Algorithm~\ref{alg:DY}, is built from the \emph{exact} change--of--measure
identity of Theorem~\ref{thm:DY}, valid for the entire admissible class
(any centered, stationary--increment Gaussian $\eta$ with continuous paths,
$\eta(0)=0$, regularly varying variance of index $\alpha\in(0,2)$; the proof uses
only \eqref{eq:cov} and the Gaussian Cameron--Martin formula). Its per--replication
output is bounded by $(n+1)/T$, giving variance uniformly controlled in $T$
(Proposition~\ref{prop:bdd}) and the distribution--free confidence bound
\eqref{eq:hoeffding}; its bias splits into an explicit discretisation part
(Lemma~\ref{lem:disc}, with bound $\Delta(T,\delta)\to0$ as $\delta\to0$ \emph{for
fixed $T$}, and quantitative rate under local self--similarity) and an explicit,
nonnegative truncation part (Lemma~\ref{lem:subadd}, $\to0$ as $T\to\infty$).
Theorem~\ref{thm:main} assembles these into mean--square consistency along the
\emph{nested} schedule ``fix $T$, then refine $\delta$, then increase $N$'', with
tuning rule $N\ge 9(1/\delta+1/T)^2/\varepsilon^2$ and $T,\delta$ chosen so each
bias term is $\le\varepsilon/3$, certifying RMSE $\le\varepsilon$. The crucial
correction to a tempting but false claim is that $\delta$ must shrink \emph{fast
enough relative to $T$} (precisely $\Delta(T,\delta)\to0$), not along an arbitrary
joint schedule, because the discretisation-bias bound grows with $T$.

\end{document}
